\section{Adding pendant edges}\label{sec:pendant}
In this section, we prove the lemma which carries out one step of the recursion sketched in Section~\ref{sec:sketch-linear}.
Given rainbow spanning trees of $K[A]$ and a set $H$ of edges from $A$ to $A'$, we give each tree one edge of $H$ at every vertex of $A'$, so that every tree remains rainbow and gains few edges with colour in any of a given family of colour sets.

\smallskip

\noindent\textbf{Notation.} Sections~\ref{sec:pendant}--\ref{sec:gap} contain the proof of Theorem~\ref{thm:gap}, and we use the following notation in them. Throughout, $\phi$ is a proper edge-colouring of a complete graph.
For a colour $\gamma$, $E_\gamma$ is the set of edges with colour $\gamma$ (a matching), and $V(\gamma)$ is the set of $2|E_\gamma|$ vertices it covers.
For a vertex $v$, $C(v)=\{\phi(vu):u\neq v\}$ is the set of colours at $v$.
For a vertex set $A$, $K[A]$ is the complete graph on $A$, and, for disjoint vertex sets $A$ and $A'$, $E(A,A')$ is the set of edges with one vertex in $A$ and one in $A'$.
For a graph $G$, $\phi(G)$ is the set of colours appearing on $E(G)$, and, for a set $Z$ of colours, $E_Z(G)$ is the set of edges of $G$ with colour in $Z$.
For a tree $T$ and a set of colours $Z$, we write $\ell_Z(T):=|\phi(T)\cap Z|$ for the number of colours of $T$ which lie in $Z$.
For integers $x\geq s\geq 0$, $\ff{x}{s}=x(x-1)\cdots(x-s+1)$, and $\log_2$ is the logarithm to base $2$.

\smallskip

Throughout this section, $\phi$ is a proper edge-colouring of a complete graph containing disjoint vertex sets $A$ and $A'$.
We first fix some notation for the lemma. Let $T_1,\dots,T_q$ be pairwise edge-disjoint rainbow spanning trees of $K[A]$, and let $H\subseteq E(A,A')$ be such that $d_H(v)=q$ for every $v\in A'$. For each $v\in A'$, let $H(v)$ be the set of edges of $H$ at $v$.
For each colour $\gamma$, let
\[
F_\gamma=\{i\in [q]:\gamma\in\phi(T_i)\},\qquad a_\gamma=|F_\gamma|\qquad\text{and}\qquad d_\gamma=|H\cap E_\gamma|,
\]
so that $a_\gamma$ is the number of trees containing colour $\gamma$ and $d_\gamma$ is the number of edges of $H$ with colour $\gamma$. Finally, let
\[
r=\max_{v\in A',\,i\in[q]}\bigl|\phi(H(v))\cap\phi(T_i)\bigr|,
\]
which is the largest number of colours that a single tree shares with the edges of $H$ at a single vertex of $A'$.
Given bijections $\lambda_v:H(v)\to[q]$, $v\in A'$, we write $\lambda(e)=\lambda_v(e)$ for each $e\in H(v)$, and, for each $i\in [q]$, we let $T_i^+:=T_i\cup\{e\in H:\lambda(e)=i\}$ be the graph formed from $T_i$ by adding the edges of $H$ given to it.

\begin{lemma}[Adding pendant edges]\label{lem:attach}
Let $q\ge4$, let $T_1,\dots,T_q$ be pairwise edge-disjoint rainbow spanning trees of $K[A]$, and let $H\subseteq E(A,A')$ satisfy $d_H(v)=q$ for every $v\in A'$; define $a_\gamma$, $d_\gamma$ and $r$ as above.
Let $M\ge1$, let $t\geq 2$ be an integer, let $\cZ$ be a family of colour sets, and suppose that the following hold for some $K$ and $\Lambda$.
\begin{props}
\item $a_\gamma+d_\gamma\le M$ for every colour $\gamma$.\label{att:mult}
\item Every colour lies in at most $K$ members of $\cZ$, and, for every $Z\in\cZ$, at most $\Lambda$ edges of $H$ have colour in $Z$.\label{att:Z}
\item $q\ge 8r+32(M-1)$, $t-1\ge 4e\Lambda/q$ and $2^{t-1}\ge16K$.\label{att:ineq}
\end{props}
Then there are bijections $\lambda_v:H(v)\to[q]$, $v\in A'$, such that each $T_i^+$ is a rainbow spanning tree of $K[A\cup A']$ in which, for each $Z\in\cZ$, fewer than $t$ of the new edges have colour in $Z$.
\end{lemma}

The quantity $r$ is the one that must be kept small throughout the recursion, as discussed in Section~\ref{sec:sketch-linear}: it appears in the first inequality in~\ref{att:ineq}, which, together with~\ref{att:mult}, ensures that the trees can be kept rainbow.
The parameter $t$ is the price we pay at each step, in that the number of colours from each $Z\in\cZ$ in each tree may increase by up to $t-1$. By the second and third inequalities in~\ref{att:ineq}, $t$ needs to be at least a constant times $\Lambda/q$, and at least logarithmic in $K$.

\begin{proof}[Proof of Lemma~\ref{lem:attach}]
Choose the bijections $\lambda_v$, $v\in A'$, independently and uniformly at random. Each graph $T_i^+$ is a spanning tree of $K[A\cup A']$, since it is formed from $T_i$ by adding one pendant edge at each $v\in A'$.
We consider the following three families of events, all of which are canonical in the sense of Section~\ref{sec:bijections}, where the coordinates are the vertices $v\in A'$ and, for each $v$, the bijection is $\lambda_v:H(v)\to [q]$.
\begin{itemize}
\item For each $v\in A'$, $e\in H(v)$ and $i\in F_{\phi(e)}$, let $A^1_{v,e,i}$ be the event that $\lambda_v(e)=i$. This event has probability $1/q$.
\item For all distinct $v,w\in A'$, all $e\in H(v)$ and $e'\in H(w)$ with $\phi(e)=\phi(e')$, and each $i\in [q]$, let $A^2_{e,e',i}$ be the event that $\lambda_v(e)=\lambda_w(e')=i$. This event has probability $q^{-2}$.
\item For each $Z\in\cZ$, $i\in[q]$ and each set $S$ of $t$ edges of $H$ with colour in $Z$ which lie at $t$ distinct vertices of $A'$, let $A^3_{Z,i,S}$ be the event that $\lambda(e)=i$ for every $e\in S$. This event has probability $q^{-t}$.
\end{itemize}
If none of these events occurs, then the conclusion of the lemma holds.
Indeed, if no event $A^1_{v,e,i}$ occurs, then no tree $T_i$ receives an edge whose colour already appears in $T_i$.
If no event $A^2_{e,e',i}$ occurs, then no tree receives two edges of the same colour at different vertices of $A'$, while two edges at the same vertex have different colours as $\phi$ is proper.
If no event $A^3_{Z,i,S}$ occurs, then, as the new edges of $T_i^+$ lie at distinct vertices of $A'$, fewer than $t$ of them have colour in $Z$.

By Lemma~\ref{lem:negdep}, the graph of adjacencies between these events is a negative dependency graph.
We give each event $A^1_{v,e,i}$ the weight $x=2/q$, each event $A^2_{e,e',i}$ the weight $x=(2/q)^2$, and each event $A^3_{Z,i,S}$ the weight $x=(2/q)^t$; all of these lie in $(0,1)$ because $q\ge4$.
Fix a term $(v,e,i)$, and let $\gamma=\phi(e)$. We bound the total weight of the events which have a term with coordinate $v$ and either domain element $e$ or value $i$, as follows.
\begin{itemize}
\item There are at most $a_\gamma$ events $A^1_{v,e,j}$, and at most $r$ events $A^1_{v,e',i}$. So the events $A^1$ contribute weight at most $2(a_\gamma+r)/q$.
\item There are at most $q(d_\gamma-1)$ events $A^2$ with a term $(v,e,j)$ for some $j$. There are at most $\sum_{e'\in H(v)}(d_{\phi(e')}-1)\le q(M-1)$ events $A^2$ with a term $(v,e',i)$ for some $e'$. So the events $A^2$ contribute weight at most $4(d_\gamma-1)/q+4(M-1)/q$.
\item For each of the two kinds of shared term, there are at most $q$ choices of the value or of the edge of $H(v)$, at most $K$ choices of $Z$ by~\ref{att:Z}, and at most $\binom{\lfloor\Lambda\rfloor}{t-1}$ choices of the remaining edges of $S$. So, by~\ref{att:ineq}, the events $A^3$ contribute weight at most
\[
2qK\Bigl(\frac{e\Lambda}{t-1}\Bigr)^{t-1}\Bigl(\frac 2q\Bigr)^t=4K\Bigl(\frac{2e\Lambda}{q(t-1)}\Bigr)^{t-1}\le4K2^{-(t-1)}\le\frac14.
\]
\end{itemize}
Since $2a_\gamma+4(d_\gamma-1)\le4(a_\gamma+d_\gamma-1)\le4(M-1)$ by~\ref{att:mult}, the events $A^1$ and $A^2$ contribute in total at most $(2r+8(M-1))/q\le 1/4$, by~\ref{att:ineq}. Thus, the total weight of these events is at most $1/2$.

An event with $s$ terms (where $s=1$, $2$ or $t$ respectively for the three families) has all of its neighbours among the events counted above for its $s$ terms. Hence, the product of $(1-x)$ over its neighbours is at least $\prod_{\text{terms}}(1-\sum x)\ge2^{-s}$, and so its weight times this product is at least $(2/q)^s2^{-s}=q^{-s}$, which equals its probability.
Therefore, by Lemma~\ref{lem:LLL}, with positive probability none of the events occurs, and the lemma follows.
\end{proof}
